% TOP_PROBLEM: {"id": 8700027, "title": "Conjecture 3.9 — Schneider", "category_id": 3, "category": {"id": 3, "name": "graph_theory", "display_name": "Graph Theory", "description": "Problems involving graphs, networks, and their properties.", "slug": "graph-theory", "order_index": 3, "created_at": "2026-07-31T15:26:25.671Z"}, "status": "open", "problem_number": "AMR-086-0027", "set_id": 13, "statusQualification": "Restores the shared nonzero-logarithm and algebraic-input hypotheses from the primary source; all powers use the same branch."}
% FIRST_POSTED: 2026-10-01
% ENTRY_KIND: statement_only
% UnsolvedMath ID 8700027; source catalogue code AMR-086-0027
% !TeX program = lualatex
% Definition and sources only; this file is not an LLM solution attempt.
\documentclass[11pt,a4paper]{article}
\usepackage[margin=25mm]{geometry}
\usepackage{fontspec}
\setmainfont{lmroman10-regular.otf}[BoldFont=lmroman10-bold.otf,
 ItalicFont=lmroman10-italic.otf,BoldItalicFont=lmroman10-bolditalic.otf]
\usepackage{amsmath,amssymb,mathtools}
\usepackage{unicode-math}
\setmathfont{latinmodern-math.otf}
\AtBeginDocument{\let\setminus\smallsetminus}
\usepackage{xurl,hyperref}
\hypersetup{colorlinks=true,urlcolor=blue}
\setcounter{secnumdepth}{0}
\setlength{\parindent}{0pt}
\setlength{\parskip}{5pt}
\setlength{\emergencystretch}{3em}
\begin{document}
\section*{Conjecture 3.9 — Schneider}\label{8700027}
{\small UnsolvedMath ID 8700027 · AMR-086-0027 · Graph Theory · AMR Open Problem Lists}
\subsection{Definitions and mathematical statement}
Let $\beta$ be an algebraic irrational complex number of degree $d=[\mathbb Q(\beta):\mathbb Q]\ge2$. Let $\alpha$ be a nonzero algebraic number, and fix a complex number $\lambda\ne0$ satisfying $e^\lambda=\alpha$. Define all powers using this same logarithm:
\[
u_j=e^{\beta^j\lambda}=\alpha^{\beta^j},\qquad 1\le j\le d-1.
\]
In particular, $\alpha^{\beta^j}$ means exponentiation by $\beta^j$, not the $j$th power of $\alpha^\beta$.

Schneider's conjecture asks whether the $d-1$ numbers $u_1,\ldots,u_{d-1}$ are algebraically independent over $\mathbb Q$. That is, must $F(u_1,\ldots,u_{d-1})$ be nonzero for every nonzero polynomial $F\in\mathbb Q[X_1,\ldots,X_{d-1}]$, for every permitted $(\alpha,\beta,\lambda)$?

These hypotheses are supplied by the original source's setup preceding the conjecture. The degree is the minimal-polynomial degree of $\beta$, not a freely chosen upper bound. A single chosen logarithm is used throughout, and it need not be the principal logarithm. Even $\alpha=1$ is permitted when a nonzero logarithm is chosen.

The exponent list stops at $d-1$ because higher powers of $\beta$ satisfy rational linear relations involving its lower powers and one. The stated conclusion excludes every rational polynomial relation among the indicated exponential values. It requires more than pairwise distinctness, pairwise algebraic independence, or transcendence of each individual entry when $d-1>1$.
\subsection{Short English statement}
Are the powers of an algebraic number indexed by successive powers of an algebraic irrational exponent jointly algebraically independent?
\subsection{Sources}
\begin{itemize}
\item[S1] ULAM.AI and the UnsolvedMath Contributors, supplied problems.json, record 8700027 (AMR-086-0027), AMR Open Problem Lists. Catalogue identity and source transcription; CC BY 4.0.
\url{https://huggingface.co/datasets/ulamai/UnsolvedMath}.
\item[S2] Waldschmidt - Open Diophantine Problems (2004), Conjecture 3.9. Original source indexed by AMR-086-0027.
\url{https://arxiv.org/abs/math/0312440}.
\item[S3] M. Waldschmidt, Open Diophantine Problems, Moscow Mathematical Journal 4 (2004), 245–305. Author PDF supplies the surrounding definitions and hypotheses.
\url{https://webusers.imj-prg.fr/~michel.waldschmidt/articles/pdf/odp.pdf}.
\end{itemize}
\end{document}
